\section{Preliminaries}
  \paragraph{Basic notation.}
    We take the natural numbers $\NN$ to include $0$. For $n \in \NN$, we
      define the set $[n] = \{0, \dots, n-1\}$. We will employ the
      convention that $N = 2^n$; thus, $NM = 2^{n+m}$.

    Functions are defined over bitstrings unless otherwise noted, and the set
      of all functions from the set of $m$-length bitsrings $\bin^m$ to the
      set of $n$-length $\bin^n$ bitstrings is denoted $\funcspace{m}{n}$.

    We use $\idmatrix[n]$ to denote the $n$-by-$n$ identity matrix, dropping
      the subscript whenever the dimension is clear from the context.

    \hhnote{From SoK: \\
      For a positive integer $n$, we let $[n]$ denote the set $\{ 1, \dots, n \}$.
      For a bit string $x\in\{0,1\}^*$, $\card{x}$ denotes its length.
      For a finite set $\Xspace$, $\card{\Xspace}$ is the cardinality of $\Xspace$.
      For a list $\Xlist$, $\Xlist[i]$ retrieves the $i$th element of the list,
      and, for an index set $\Ispace$, $\Xlist[\Ispace]$ indicates
      the list $\Xlist$ restricted to the elements whose indices are
      contained in $\Ispace$. 
    }

    \hhnote{
      todo: introduce condprobsample
    }
